\phantomsection
\addcontentsline{toc}{chapter}{Rubrics}
\fancyhead[C]{{\LARGE Rubrics}}
\lett{I}{n} the Christian life, the Holy Sacrifice of the Mass and the Daily Office serve as the highest modes of prayer. Due to their solemnity, it is important to attend carefully to their proper celebration.

	Through the Mass, the highest sacrifice is offered unto the Most Holy Trinity, and Christ is made substantially present upon the Altar. In the Office, the Christian obediently, faithfully, and joyfully observes the Sacred Scriptures' exhortation to offer the morning and evening sacrifice.
	
	The Christian who prays this Book, both layman and cleric, must keep the rubrics and their purpose within these prayers in mind. At the same time, however, the rubrics are to help the Christian offer his Sacrifice with loving care and faithfulness, not to bog him down with obligations. The Rites and Rubrics in this Book are laid out in a manner that should be accessible, especially after a few uses, lest `there be more business to find out what should be read, than to read it when it be found out.'
%It is in this mode of prayer that the Christian obediently, faithfully, and joyfully observes the Sacred Scripture's exhortation to offer the morning and evening sacrifice. And, through devotion to the Minor Hours, to `pray without ceasing'.\par
%It must be remembered, before a review of the rubrics, that the purpose of the Office is to facilitate this Sacrifice. Therefore, a simpler Office or an Office devoutly offered with imprecise rubrical distinctions, is immensely more pleasing than a despondent spirit.\par
%\emph{
%}
\section{General Rubrics}
\begin{description}
\item[Advent Ferial Commemoration] During Advent, the Feria is commemorated with the Prayers of the First Sunday of Advent.
\item[Alleluia] The `Alleluia' is wholly omitted absolutely during Septuagesimatide and Lent.
\item[Amen] When the word \emph{Amen} is in italics, it should be said by the Congregation in response to the Minister's prayer.
%This modifies the 1928 rubric, in line with the 1549 and 1662, especially with consideration of Orthodox doctrine.
\item[Antecommunion] It may occur that a community lacks a Priest on Sunday or another day where the Service of Holy Communion is desired. In such a case, Mattins and the Great Litany should first be said. Then, the Holy Mass may be said from the \emph{Collect for Purity} until the Sermon, inclusive, concluding with the Additional Collects (p. \pageref{AdditionalCollects}).\par
\textsc{Note,} Mattins, the Great Litany, and Antecommunion should all be led by the Minister in the same place as the Office, not at the Altar.
%From the Monastic Diurnal
\item[Collects] When several proper Collects (or Secrets, or Postcommunions) are said, only the first and the last receive the full ending.\par
	\textsc{Note,} In the Mass, the \emph{Lord be with you} is only said to introduce the first and second collects.
\item[Commemorations] When multiple Feast Days overlap on the same day, the Office will be of the Feast Day of highest rank. The other Feast Day(s) will then be commemorated (or, in some cases, suppressed or transferred). Commemoration consists of praying the highest-ranking Feast Day's collect first followed by the lower ranking Feast Day(s)' collect. In the Mass, this also applies to the Secret and Postcommunion.

\textsc{Note,} Once three Collects (Secrets, Postcommunions) have been said, further Commemorations cease and are omitted that year.
%Collect Conclusion taken from the English Missal 1940 (it's not provided in the Lay Missal).
%\item[Conclusion of Collects] When the Collect is addressed to God the Father, the ending is `Through Jesus Christ thy Son our Lord; who liveth and reigneth with thee in the unity of the Holy Spirit, one God, world without end. Amen.'\par
%If in the beginning of the Collect the Son be mentioned: `Through the same Jesus Christ thy Son our Lord; who liveth and reigneth with thee in the unity of the Holy Spirit, one God, world without end. Amen.'\par
%If at the end of the Collect the Son be mentioned: `Who with thee liveth and reigneth in the unity of the Holy Spirit, one God, world without end. Amen.'\par
%If the Collect be addressed to the Son: `Who livest and reignest with God the Father in the unity of the Holy Spirit, one God, world without end. Amen.'\par
%When the Holy Spirit is mentioned in the Collect: `In the unity of the same Holy Spirit, etc.'
%Taking what I can from the 1933 Lay Missal:
\item[Conclusion of Collects] When the Collect is addressed to God the Father, the ending is `Through Jesus Christ thy Son our Lord, who liveth and reigneth with thee, in the unity of the Holy Ghost, God, throughout all ages, world without end. Amen.'\par
If in the beginning of the Collect the Son be mentioned: `Through the same Jesus Christ thy Son our Lord; who liveth and reigneth with thee in the unity of the Holy Ghost, God, throughout all ages, world without end. Amen.'\par
If at the end of the Collect the Son be mentioned: `Who liveth and reigneth with thee in the unity of the Holy Ghost, God, throughout all ages, world without end. Amen.'\par
If the Collect be addressed to the Son: `Who livest and reignest with the Father, in the unity of the Holy Ghost, God, throughout all ages, world without end. Amen.'\par
If the Collect be addressed to the Holy Spirit: `Who livest and reignest with the Father and the Son, God, throughout all ages, world without end. Amen.'\par
When the Holy Spirit is mentioned in the Collect: `In the unity of the same Holy Ghost, etc.'
\item[Dominus Vobiscum] The Minister may only proclaim the \emph{Dominus vobiscum} (The Lord be with you) if he be in Major Orders. Otherwise, he must instead lead with the \emph{Domine, exaudi orationem nostram} (O Lord, hear our prayer) wherever the \emph{Dominus vobiscum} appears, or omit it entirely if it would be said twice consecutively.
\item[Ferial Days] The Ferial Office, that is, the Simple Office of the occurrent Season, is always said on weekdays in Lent, on Ember Days, and on Rogation Monday, unless a I or II Double occur, in which case the Office is of the Feast with Commemoration of the Feria.\par
On weekdays in Advent and between Septuagesima and Ash Wednesday, and on Common Vigils, the Ferial Office is always said unless a Double Feast or Octave occur, and then Commemoration is made of such Ferias.\par
If a Simple Octave Day occur on one of these Ferias, it is only commemorated.\par
In like manner throughout the year the Office is of the Feria on those weekdays on which there does not occur a Double Feast, an Octave, or the Office of St. Mary on Saturday.
\item[Gloria Patri] During Passiontide, the \emph{Gloria Patri} is omitted in the Invitatory of Mattins, Responsory of the Minor Hours, Aspersion of the People, Introit, and Lavabo. During the Sacred Triduum, the \emph{Gloria Patri} is wholly omitted.
\item[Latin] In much of this Book, Latin is provided alongside the English. The Latin may always be used instead of the English, unless otherwise determined by the Ordinary. When Latin is not provided, but great devotion exists for the Latin language, the Latin which serves as the base text (such as the Liber Precum Publicarum or Breviarium Monasticum) may be used, with great concern for the catechesis and participation of the People.
\item[Lenten Ferias] When the Lenten Feria is indicated, the propers of Ash Wednesday may be said.
\item[Liturgical Seasons] For the purpose of the rubrics, these festive seasons are referenced:
    \begin{itemize}
    	\item Advent: From I Evensong of the First Sunday of Advent until I Evensong of Christmas Day, exclusive.
        \item Christmastide: From I Evensong of Christmas Day until I Evensong of Epiphany Day, exclusive.%\footnote{This is not to be confused with the `Christmas season', which can be said to end on Candlemas.}
        \item Epiphanytide: From I Evensong of Epiphany Day until I Evensong of Septuagesima Sunday, exclusive.
        \item Septuagesimatide (Pre-Lent): From I Evensong of Septuagesima Sunday until Mattins of Ash Wednesday, exclusive.
        \item Lent: From Mattins of Ash Wednesday until I Evensong of Passion Sunday, exclusive.
        \item Passiontide: From I Evensong of Passion Sunday until Mattins of Maundy Thursday, exclusive.
        \item Sacred Triduum: From Mattins of Maundy Thursday until I Evensong of Easter Sunday, exclusive.
        \item Eastertide (Paschaltide): From I Evensong of Easter Sunday until I Evensong of Trinity Sunday, exclusive.
        \item Trinitytide: From I Evensong of Trinity Sunday until I Evensong of the First Sunday of Advent, exclusive.
    \end{itemize}
\item[Minor Feast Days] Since only I and II Doubles are provided here in this Book (with some exceptions), the rubrics will be geared towards those. When the Office or Mass is celebrated with the rest of the Feast Days, the rubrics of either the \emph{Saint Botolph's Breviary} or \textit{Saint Botolph's Missal} should be consulted.%\par
%The Mass propers and Office antiphons for the Feast Days not supplied in this Book may be found in the \textit{Book of Occasional Services}.\par
%For the rest of the Office propers, the Monastic Breviary should be consulted.

%Is this commentary really appropriate for the rubrics?:
% However, since these books have been published by Anglicans, great care must be taken to review the text in order to avoid any heterodox prayers.
\item[Mode of Recitation] Wherever it is indicated something be read, it may also be chanted, except for the Fore-Office and the Prayers after the Third Collect.
\item[Octaves] The Office of an Octave is said (or Commemoration thereof made, when it is hindered by a Feast or a Sunday) through eight continuous days in Octaves of Feasts of the I Class. Octaves of Feasts of the II Class, which are Simple Octaves, are kept only on the Octave Day itself, with the rite of a Simple, unless hindered by a more worthy Office; but no notice is taken of the days within the Octave. If any Octave is not ended before Ash Wednesday, the Vigil of Pentecost, or 17 December, no notice is taken of it thenceforth.\par
A Simple Octave Day occurring within an Octave is only commemorated.
\item[Octave Days] On the Octave Day, except for the proper Psalms and Lessons, the Office is said as on the day of the Feast, unless otherwise noted. On a Simple Octave Day the Office is Simple.
\par
On Sundays within Privileged Octaves, the Office is said as directed in the Proper of Season. On Sundays within Common Octaves, the Office is of the Sunday with Commemoration of the Octave.
\item[Office of Our Lady on Saturday] On all Saturdays---except in Advent, Septuagesimatide, Lent, and on Ember Days---unless the Office be of a Double Feast (even transferred), or of an occurrent Octave or Vigil, or of a Sunday anticipated according to the rubrics, the Office is of St. Mary.
%\item[Plural \& Gendered Nouns and Pronouns] Sometimes there are words in the liturgy which change depending on number and gender, such as if the person being prayed for is a woman or if there are multiple people being prayed for. The Book, in accordance with English grammar, assumes the masculine singular, and if the prayer should be changed for either number or gender, the relevant words and letters are italicised.\par
%In cases where the change could only be due to number, such as prayers for Clerics (Subdeacons, Deacons, Priests, Bishops, etc.) which are---by divine law---always masculine, the relevant words are still italicised.
\item[Sunday Collects] %The Collect for a Sunday is said on the I Evensong, Mattins, and II Evensong of the Sunday, unless a Feast Day of higher rank is celebrated. It is also said throughout the week, unless a proper Collect is provided for the weekday.
If the Office be of a Feria without a proper Collect, then the Sunday Collect is read. Otherwise, the Collect for the Feast Day is read instead.
\end{description}

\section{Office Rubrics}
\begin{description}
%\item[Canticle Antiphons] %The antiphons for the \emph{Benedictus} and \emph{Magnificat} canticles are provided for those with special devotion to them.\par
%On Feast Days below Double, the antiphon may be said: as far as the {\dag} before, and repeated in full after, each Canticle.\par
%On Feast Days Double and higher, the antiphon may be said: entire before and after each Canticle.%\par
\item[Collects] The Collect to be said during the Office is provided in the relevant Proper. The Seasonal Prayers are to be said during Mass as indicated. The Seasonal Prayers (only its Collects) may be said as additional Prayers, as indicated in the Daily Office.
%The antiphons, in line with the Prayer Book Tradition, are optional.
\item[Gloria Patri] After every Psalm and the indicated Canticles, except in the Office of the Dead or where otherwise indicated, the \emph{Gloria Patri} should be said.\par
It is customary to reverence the Trinity by bowing one's head during the \emph{Gloria Patri}, instead of crossing oneself.
%\item[Great Litany] The petition `From all sedition, privy conspiracy, and rebellion' may be replaced by the version from the 1549 Book of Common Prayer.
\item[Lectionary] The Lessons and Psalms for the days of the Church Year are provided in the lectionary, according to the Rankings of Feast Days (p. \pageref{ranking}).\par
\textsc{Note,} The proper psalms for Second Class and lower Sundays---as well as Ember Days---are optional.

%\item[Octave Day Lessons] Only on the Feast Day are the proper Psalms and Lessons of the Feast Day used, not also on the Octave Day. During the Octave, the Collect for the Feast is used, according to relevant rules for Octave rankings.
\item[Office of the Dead] The Office of the Dead may be prayed in addition to the Daily Office. It may also be prayed instead of the Daily Office on an ordinary Feria, a Memorial, or a Feast of Simple rank.\par
%CHECK:
The Office of the Dead may also replace the Daily Office of rank Double or lower on the day of death and the day of burial.
\item[Office of the Dead Antiphons] The Antiphons are not doubled except on 2 November; %14 November (at the Commemoration of All the Departed O.S.B., for Benedictines);
on the day of a burial; on the day after receiving tidings of a death; on the third, seventh, and thirtieth days; on the anniversary (even when transferred); and whenever the Office is celebrated solemnly.
\item[Office of the Dead Lessons] The lessons are always read without introduction or conclusion.
%\item[Proper Psalms \& Lesson] Proper Lessons are given for every day of the Church Year and for every Feast Day supplied in this Book (Prayer Book Holydays). Proper Psalms are supplied for every Sunday, Ember Days, and Prayer Book Holydays.\par
%Only Feast Days of II Double or higher receive Proper Psalms and Lessons, unless otherwise noted or supplied in this Book.
\item[Psalm Antiphons] The Psalm Antiphons may be used at the discretion of the Minister.
\item[Scripture Verse Numbers] When a lesson or psalm is listed with only one verse (such as John 1:45 or Psalm 102:15), the Minister reads the that verse and the verses following for the rest of the chapter.
\item[Suffrages] The Suffrages in the Minor Hours are not said on Feast Days Double or higher.
\item[Third Lesson] In accordance with pious custom, after the Second Canticle in the Major Hours, but before the Creed, it is fitting for a patristic lesson to be read from the Feast Day's Second or Third Nocturn in Monastic or Sarum Matins.
\item[Trinitytide Opening Sentences] The Opening Sentences for Trinity Sunday may also be said on the Sundays after Trinity.
\item[Vigils] %Vigils, unlike modern Roman practice, are not anticipated celebrations of the following day. Rather, they are days in their own right, dedicated to fasting in preparation for the Feast Day.\par
The Office of a Vigil is only said at Mattins.\par
If a Vigil fall on a Sunday (except for the Vigils of Christmas \& Epiphany), it is said or commemorated on the preceding Saturday.
\end{description}